\subsection{缩塌子空间所得的空间}

设 X 为拓扑空间，A 为 X 的子集。考虑 X 上最细的等价关系 R，使 A 的所有元素彼此等价：X 的两个元素关于这个关系等价，当且仅当它们相等，或者它们都属于 A。

定义 10. --- X 对 R 的商空间记作 X/A，称为由 X 缩塌子集 A 所得的空间。

用 \(\rho\colon X\to X/A\) 表示典范满射。设 Y 为拓扑空间，\(f\colon X\to Y\) 为连续映射。要存在连续映射 \(g\colon X/A\to Y\) 满足 \(g\circ\rho=f\)，必须而且只需 \(f\) 在 A 上是常值的。

若集合 A 在 X 中闭（相应地，开），则映射 \(\rho\) 是闭（相应地，开）映射，并诱导出 \(X\setminus A\) 到其像上的一个同胚。若 A 是 X 的闭且拟紧的子集，则映射 \(\rho\) 是逆紧的（TG, I, p. 75, 定理 1）。

若集合 A 为空，则 \(\rho\) 是同胚。若 A 非空，则 A 缩塌为 X/A 的一点，该点称为 X/A 的基点，记作 \(s_{\mathrm{X / A}}\)。这时把空间 X/A 等同于沿从 A 到 \(\{s_{\mathrm{X / A}}\}\) 的唯一映射把空间 X 附着到空间 \(\{s_{\mathrm{X / A}}\}\) 上所得的空间。

命题 10. --- 设 X 为拓扑空间，A 为 X 的子集。假设存在同伦 \(\sigma: \mathrm{X} \times \mathbf{I} \to \mathrm{X}\)，以 \(\mathrm{Id}_{\mathrm{X}}\) 为初始，在 \(\mathrm{A} \times \{1\}\) 上是常值的，且满足 \(\sigma(\mathrm{A} \times \mathbf{I}) \subset \mathrm{A}\)。那么从 X 到 X/A 的典范映射 \(\rho\) 是同伦等价。

用 \(f\) 表示 \(\sigma\) 的末项。它是从 X 到 X 中的连续映射，同伦于 \(\mathrm{Id}_{\mathrm{X}}\)。它在 A 上是常值的；因此存在唯一的连续映射 \(g \colon \mathrm{X} / \mathrm{A} \to \mathrm{X}\)，满足 \(g \circ \rho = f\)。另一方面，由于 \(\sigma(\mathrm{A} \times \mathbf{I}) \subset \mathrm{A}\)，存在映射 \(\sigma'\)，它从 \(\mathrm{X} / \mathrm{A} \times \mathbf{I}\) 到 \(\mathrm{X} / \mathrm{A}\) 中，满足 \(\sigma'(\rho(x), t) = \rho(\sigma(x, t))\)（对所有 \(x \in \mathrm{X}\) 及每个 \(t \in \mathbf{I}\)）。由 I, p. 19, 命题 10，映射 \(\sigma'\) 是连续的。它是连接 \(\mathrm{X} / \mathrm{A}\) 的恒同映射与映射 \(\rho \circ g\) 的同伦。因此，映射 \(g\) 与 \(\rho\) 是互为同伦逆的同伦等价。

注记。--- 设 X 为拓扑空间，A 为 X 的子集。

\begin{enumerate}
\def\labelenumi{\arabic{enumi})}
\item
  若从 X 到 X/A 上的典范映射是同伦等价，则存在同伦 \(\sigma\colon\mathrm{X}\times\mathbf{I}\to\mathrm{X}\)，它以 \(\mathrm{Id}_{\mathrm{X}}\) 为初始，且在 \(\mathrm{A}\times\{1\}\) 上为常值。但也可能不存在除此之外还满足条件 \(\sigma(\mathrm{A}\times\mathbf{I})\subset\mathrm{A}\) 的同伦（III, p. 322, 习题 4）。
\item
  可能存在连接 X 的恒同映射与 X 到 X 中、在 A 上为常值的映射的同伦，而空间 X 与 X/A 却不同伦等价（III, p. 325, 习题 14）。
\item
  假设 A 是可缩的。若偶 (X, A) 具有同伦扩张性质，则从 X 到 X/A 的典范映射是同伦等价。事实上，设 \(\sigma\colon A \times I \to A\) 为一同伦，其初始映射是 A 的恒同映射，末项映射是常值映射。于是存在同伦 \(\sigma'\)，它以 \(\mathrm{Id}_{X}\) 为初始，且是 \(\sigma\) 的开拓。断言便由命题 10 得出。
\item
  设 X 与 Y 为拓扑空间，A 为 X 的非空子空间，B 为 Y 的子空间。设 \(f: X \to Y\) 为连续映射，满足 \(f(A) \subset B\)。用 \(\varphi: X/A \to Y/B\) 表示由 f 在商上诱导的连续映射。若 f 定义了偶 (X, A) 到偶 (Y, B) 上的同伦等价，则映射 \(\varphi\) 是偶 (X/A, \(s_{X/A}\)) 到偶 (Y/B, \(s_{Y/B}\)) 上的同伦等价。设 P 为缩化为单独一点的拓扑空间，\(\alpha\) 与 \(\beta\) 为从 A 与 B 到 P 中的常值映射。注意到映射 \(\varphi\) 可等同于由 f 与 P 的恒同映射诱导的、从 \(P \cup_{\alpha} X\) 到 \(P \cup_{\beta} Y\) 中的映射。断言便由 III, p. 249 的推论 1 得出。
\end{enumerate}

